\chapter{Introduction}

Insiders threats have become one of the most significant issues in cybersecurity nowadays. Differently from external hackers, they are insiders, i.e., legitimate users that have authorized access to organizational systems, so they can use their privileges in a wrong way either for a purpose of a kind of malicious or for that of a negligent \cite{bishop2008defining}. Since they know the internal processes as well as the system architecture, they can dodge usual security measures, thus, they are a source of undetected data leakages, financial losses, and operational disruptions in most cases \cite{inayat2024insider}.

Industry reports have been, gradually, over the last ten years, drawing attention toward a significant increase in the frequency and cost of insider incidents \cite{worth2025databreach}. The Ponemon Institute’s 2022 report \cite{Ref3} showed an increase of 44\% in such events as compared to 2020, and the average organizational loss went up from \$11.45 million to \$15.4 million. The survey conducted by Gurucul in 2023 \cite{gurucul2024motivationalmisuse} found that 74\% of companies had been victims of insider attacks, while the IBM Security 2024 research \cite{IBM_2024}, which surveyed 83\% of organizations, reported insider-related incidents. The average cost per incident was estimated to be \$15.38 million. Moreover, the Ponemon study 2025 \cite{dtex2025insiderrisks,verizon2025dbir-emea} unveiled that the annual losses had escalated to \$17.4 million, thus pointing out that the detection of insiders through the currently employed methods is mostly unsuccessful and there is a need to implement adaptive and interpretable methods urgently.

Insider threats are classified into two groups: intentional and unintentional ones \cite{hunker2011insiders}. Intentional threats are the results of a malicious insider’s actions such as disgruntled employees or external collaborators who copy data, commit fraud, or cause systems to fail. While unintentional threats can be caused by human mistakes, misconfigured systems, or people who are easily influenced by phishing. However, even though awareness programs and policy enforcement can significantly reduce the rate of unintentional ransomware, it is still very difficult to identify the network of intentional ones because the latter behave covertly and their activities closely resemble legitimate user ones.

While advances in Machine Learning (ML) and Deep Learning (DL) have strengthened intrusion detection systems, insider threat detection continues to face significant limitations \cite{alzaabi2024review,gonzales2024towards}. The first challenge arises from the extreme class imbalance caused by the scarcity of malicious insider examples, which skews model performance towards majority classes. The further challenge lies in the subtlety of behavioral changes, as malicious actions often appear similar to legitimate user activities. The next challenge is the restricted interpretability of complex models, which hinders their acceptance in security-sensitive environments where decision transparency is essential. The other challenge is the dynamic drift in user behavior to execute malicious activity. These challenges lead to frequent misclassifications, high false alarm rates, and limited usefulness of current methods in real-world situations.
In addition, the growing interconnection of organizational networks with Internet of Things (IoT) devices has significantly increased the possible areas to be attacked of the most vital systems. As these devices usually operate as endpoints in larger infrastructures like smart cities and industrial control systems, they are exposed to attacks originating from the insiders or those with the help of the device. The IoT security ensuring device integrity, data confidentiality, and system availability. Having a well-planned defense strategy that unites human-centric insider detection to secure modern cyber–physical ecosystems \cite{amiri2023federated, 11082266}.

This method attempts to offer an adaptive and interpretable framework for detecting insider threats by integrating causal inference with Liquid Neural Networks (LNNs) to address the identified limitations. The suggested system uses causal inference for reasoning relationship with features to labels, and MAD to refine the selection of features. The Synthetic Minority Oversampling Technique (SMOTE) is applied to training data for oversampling the minority class to deal with the imbalance. LNNs allow the system to monitor the behavioral context changes.


The key contributions of this study are in the following order:
\begin{enumerate}
\item Evaluation of feature granularity at the session, day, and behavioral preprocessing are monitor user activities at different temporal levels.   
\item Integration of causal inference with Median Absolute Deviation (MAD) refinement to choose features that are causally significant and stable, which makes them easier to understand and reduces the number of false alarms.  
\item Addressing class imbalance through SMOTE for detecting minority class.
\item Utilizing Liquid Neural Networks (LNNs) for predicting evolving insider behaviors from time to time.  
\item A comprehensive experimental evaluation of the CMU-CERT datasets (r4.2, r5.2, r6.2), which includes domain-specific, cross-dataset, and distinctive-user testing, validated through ablation, scalability, and statistical methods.   
\end{enumerate}

 The rest of the article is arranged like this:   The review of the existing literature is in Section~\ref{literature}.   Section~\ref{background} talks about the mathematics behind the methods used in this article.   Section~\ref{proposed} gives a full explanation of the methodology.   Section~\ref{setup} talks about the setup for the experiment and metrics used.  Section~\ref{results} talks about the results analysis and the discussion.   Finally, Section~\ref{conclusion} wraps up the paper with ideas for future research and a list of sources.
